\sscp{\tsc{Atamanu}}{10-03-02}
\ili{Yalahatan} is currently spoken in a couple of villages in south central Seram.
\ili{Atamanu} is spoken in two villages on the northeastern coast of
the Elpaputih Bay (see \frf{07-03}).
Note that although \ili{Yalahatan}
and \ili{Atamanu} are said to be linguistically related to \ili{Alune},
they are separated by some distance.
\citet[38]{Collins1983} says, “The dispersal of the \ili{Atamanu} language
(Haruru, \ili{Awaiya} and \ili{Yalahatan}) may be connected with the seventeenth
and eighteenth century colonial expeditions of attrition; at
l\ili{east} this is implied in local oral traditions.”

\ili{Yalahatan} has occasional glottal-insertion between vowels (for
example: \mbox{*layaR} > \it{leʔal-e} ‘sail (n)’; *dahun > \it{laʔin
{\tl} lain} ‘leaf’; *zauq > \it{a-raʔu-e} ‘far’).
Sometimes the [ʔ] is a retention of PMP *k,
although often *k > {\0} as in *ikan ‘fish’ > \it{ian-e}.
For example: *buku > \it{huʔu-ti} ‘joint’; *paniki > \it{iniʔe}
‘fruitbat’; *bakaq > \it{haʔa} ‘split in half’; *kakak > \it{ʔaʔa-i}
‘elder sibling’; PCM **qaqi,
**kaki > \it{ʔaʔi-n} ‘foot,
leg’; *kutu > \it{ʔutu} ‘headlice, flea, tick’, etc.

While it appears that there is a merger of PMP *R/*l/*j,
the conditioned reflexes of each are different.
The usual reflex of PMP *l is \it{l,} as shown in \qf{10-39}.
Adjacent to the high vowel \it{i} there appears to be a split
of *l > \it{l, r}.
Reflexes of *taliŋa ‘ear’
and **kabil ‘fish hook’ have *l > \it{r} as shown in \qf{10-40},
while reflexes of *lima ‘five’,
*bəli ‘buy’, and **malip ‘laugh’ have *l = \it{l}.
\\

\noindent{\wg\st{6pt}\begin{tabularx}
{\linewidth}{>{\hspace{-\tabcolsep}}l<{\hspace{-\tabcolsep}}
>{\hspace{-\tabcolsep}\enspace}lll
>{\raggedright\arraybackslash}X}
\lsptoprule
%\tbx{10-39}	&	PMP/PCM	&	\ili{Yalahatan}	&	\ili{Atamanu}	&	gloss	\\	\midrule
%\endfirsthead
	&	PMP/PCM	&	\ili{Yalahatan}	&	\ili{Atamanu}	&	gloss	\\	\midrule
%\endhead
\tbx{10-39}	&	*\tbr{l}u\tbr{l}un	&	\tbr{l}u\tbr{l}u-e	&		&	roll up	\\
	&	*\tbr{l}ayaR	&	\tbr{l}eʔal-e	&	\tbr{l}ial-e	&	sail (n)	\\
	&	*\tbr{l}ima	&	\tbr{l}ima	&	\tbr{l}ima	&	five	\\
	&	*\tbr{l}aŋit	&	\tbr{l}ant-e	&	\tbr{l}anit-e	&	sky	\\
	&	*ku\tbr{l}it	&	\tbr{l}uti-n	&	\tbr{l}uti	&	skin (metathesis also in several other languages)	\\
	&	*bu\tbr{l}an	&	hu\tbr{l}an-e	&	fu\tbr{l}an-e	&	moon	\\
	&	*bu\tbr{l}u	&	hu\tbr{l}u	&	fu\tbr{l}u-i	&	body hair	\\
	&	*qu\tbr{l}u	&	u\tbr{l}u-i	&	u\tbr{l}u	&	head	\\
	&	*u\tbr{l}əj	&	u\tbr{l}e	&	(finsora)	&	caterpillar	\\
	&	*tə\tbr{l}u	&	te\tbr{l}u	&	te\tbr{l}u	&	three	\\
	&	*wa\tbr{l}u	&	wa\tbr{l}u	&	wa\tbr{l}u	&	eight	\\
	&	*qatə\tbr{l}uR	&	te\tbr{l}ul-i	&	a\tbr{l}u-i	&	egg (\ili{Atamanu} possibly non-cognate)	\\
	&	*bə\tbr{l}i	&	a-he\tbr{l}i-e	&	a-fe\tbr{l}i	&	sell (lit: cause to buy it)	\\
	&	**ma\tbr{l}ip	&	ma{\tl}ma\tbr{l}i	&	ma{\tl}ma\tbr{l}i	&	laugh	\\
	&	**ma\tbr{l}abaw	&	ma\tbr{l}aha	&	ma\tbr{l}afa	&	rat, mouse	\\
	&	**\tbr{l}əsi	&	\tbr{l}esi 	&		&	more, more than	\\
	&	**batə\tbr{l}u	&	hete\tbr{l}u	&	fete\tbr{l}u	&	thick (regional metathesis from PMP *təbəl)	\\	\midrule
\tbx{10-40}	&	*ta\tbr{l}iŋa	&	ti\tbr{r}ina-i	&	te\tbr{r}ina	&	ear	\\
	&	**kabi\tbr{l}	&	ahi\tbr{r}-e	&		&	fish hook	\\
\lspbottomrule
\end{tabularx}\mbox{}\\}

\newpage
In most cases *R > \it{l} as shown in \qf{10-41}.
A small number of etyma also show *R > \it{r},
as shown in \qf{10-42}.
*Rusuk has *R > \it{r} in \ili{Yalahatan},
but *R > \it{l} in \ili{Atamanu}.
Note the likely progressive assimilation of *R > \it{r} /rV{\gap}
in reflexes of *duRi
and *mantəR.
\\

\noindent{\wg\st{6pt}\begin{tabularx}
{\linewidth}{>{\hspace{-\tabcolsep}}l<{\hspace{-\tabcolsep}}
>{\hspace{-\tabcolsep}\enspace}lll
>{\raggedright\arraybackslash}X}
\lsptoprule
%\tbx{10-41}	&	PMP	&	\ili{Yalahatan}	&	\ili{Atamanu}	&	gloss	\\	\midrule
%\endfirsthead
	&	PMP	&	\ili{Yalahatan}	&	\ili{Atamanu}	&	gloss	\\	\midrule
%\endhead
\tbx{10-41}	&	*\tbr{R}umaq	&	\tbr{l}uma	&	\tbr{l}uma	&	house	\\
	&	*\tbr{R}ibu	&	\tbr{l}ihun	&	\tbr{l}ihun-e	&	thousand (*R > \it{l} /\it{i})	\\
	&	*wa\tbr{R}əj	&	wa\tbr{l}e-te	&	wa\tbr{l}e-te	&	vine, rope	\\
	&	*haba\tbr{R}at	&	ha\tbr{l}at-e	&	fa\tbr{l}at-a	&	west monsoon	\\
	&	*qaba\tbr{R}a	&	ha\tbr{l}a-i	&	fa\tbr{l}a	&	forearm, hand (PMP shoulder)	\\
	&	*u\tbr{R}at	&	u\tbr{l}at-a	&	u\tbr{l}at-e	&	tendon	\\
	&	*baqə\tbr{R}u	&	he\tbr{l}u	&	fe\tbr{l}u	&	new	\\
	&	*niu\tbr{R} > **niwə\tbr{R}	&	nie\tbr{l}-e	&	nie\tbr{l}-e	&	coconut	\\
	&	*qatəlu\tbr{R}	&	telu\tbr{l}-i	&	(alu-i)	&	egg (\ili{Atamanu} \it{alu-i} probably non-cognate)	\\
	&	*\tbr{R}umbia	&	\tbr{l}ipia	&	\tbr{l}ipia	&	sago	\\
	&	*ma-bə\tbr{R}əqat	&	ai-pi\tbr{l}a	&		&	heavy	\\
	&	*ku\tbr{R}ita	&	u\tbr{l}ita	&		&	octopus  (*R > \it{l} /\it{i})	\\
	&	*laya\tbr{R}	&	leʔa\tbr{l}-e	&	lia\tbr{l}-e	&	sail (n)	\\
	&	*wahiyə\tbr{R}	&	wae\tbr{l}-e	&	wae\tbr{l}-e	&	water	\\
	&	*timu\tbr{R}	&	timi\tbr{l}-e	&	timu\tbr{l}-e	&	\ili{east} monsoon  (*R > \it{l} /\it{i})	\\	\midrule
\tbx{10-42}	&	*\tbr{R}usuk	&	\tbr{r}usu ala-i	&	\tbr{l}usu	&	rib	\\
	&	*du\tbr{R}i	&	ru\tbr{r}i-n	&	ru\tbr{r}i	&	bone (assimilation)	\\
	&	*mantə\tbr{R}	&	mare\tbr{r}-e	&		&	cuscus, \it{Phalanger} (assimilation after *nt > \it{r})	\\
	&	*ma\tbr{R}uqanay	&	manuwa-i	&	manuwa-i	&	male	\\
\lspbottomrule
\end{tabularx}\mbox{}\\}

While *maRuqanay has *R > {\0},
this also occurs in several \tsc{\ili{Ambon}-Seram} languages for this item.
Examples include: \ili{Paulohi}, \ili{Sepa},
\ili{Nusalaut}, \ili{Saparua}, and \ili{Wemale} \it{manawa}, as well as \ili{Alune} \it{manukwai {\tl} makwai}.

PMP *j > \it{l},
as shown in \qf{10-43}.
Note that final *j > {\0} in the two reflexes which have been identified.
This is different to final *R/*l > \it{l},
seen in \qf{10-39}--\qf{10-41} above,
and is potentially significant in showing that \tsc{\ili{Ambon}-Seram} *j/*l
> **l can only be posited for non-final position.

\noindent{\wg\st{6pt}\tblr{>{\hspace{-\tabcolsep}}l<{\hspace{-\tabcolsep}}
>{\hspace{-\tabcolsep}\enspace}
llll}
\lsptoprule
%\tbx{10-43}	&	PMP	&	\ili{Yalahatan}	&	\ili{Atamanu}	&	gloss	\\	\midrule
%\endfirsthead
\tbx{10-43}	&	PMP	&	\ili{Yalahatan}	&	\ili{Atamanu}	&	gloss	\\	\midrule
%\endhead
	&	*pa\tbr{j}ay	&	ha\tbr{l}a	&	fa\tbr{l}a	&	‘rice’	\\
	&	*ŋa\tbr{j}an	&	na\tbr{l}a-i	&	na\tbr{l}a	&	‘name	\\
	&	*hua\tbr{j}i	&	wa\tbr{l}i-n	&	wa\tbr{l}i	&	‘younger sibling (same sex)’	\\
	&	*ma\tbr{j}a	&	ma\tbr{l}a	&		&	‘dry up’	\\
	&	*qapə\tbr{j}u	&	o\tbr{l}u-i	&	o\tbr{l}u	&	‘gall, bile’ (irregular *ə > \it{o})	\\
	&	*ulə\tbr{j}	&	ule	&	(finsora)	&	‘caterpillar’ (*j > {\0} /{\gap}{\#})	\\
	&	*waRə\tbr{j}	&	wale-te	&	wale-te	&	‘vine, rope’ (*j > {\0} /{\gap}{\#})	\\
\lspbottomrule
\end{tabular}\mbox{}\\}

While there is almost a complete merger of *R/*l/*j > \it{l}
in \ili{Yalahatan}, the different reflexes of *j word finally show that the merger
is not complete in all environments.

\largerpage
\ili{Yalahatan} has *z > \it{r},
as shown in \qf{10-44}, and *d > \it{r,
l} as shown in examples \qf{10-45} and \qf{10-46} respectively.
\ili{Atamanu} is more regular with *z/*d > \it{r} in all identified reflexes.
\\

\noindent{\wg\st{6pt}\begin{tabularx}
{\linewidth}{>{\hspace{-\tabcolsep}}l<{\hspace{-\tabcolsep}}
>{\hspace{-\tabcolsep}\enspace}lll
>{\raggedright\arraybackslash}X}
\lsptoprule
%\tbx{10-44}	&	PMP	&	\ili{Yalahatan}	&	\ili{Atamanu}	&	gloss	\\	\midrule
%\endfirsthead
	&	PMP	&	\ili{Yalahatan}	&	\ili{Atamanu}	&	gloss	\\	\midrule
%\endhead
\tbx{10-44}	&	*\tbr{z}auq	&	a-\tbr{r}aʔu-e	&	a-\tbr{r}aʔu-e	&	far	\\
	&	*\tbr{z}alan	&	\tbr{r}aran-e	&	\tbr{r}aran-e	&	path (*l > \it{r} through progressive assimilation)	\\
	&	*qu\tbr{z}an	&	u\tbr{r}an-e	&	u\tbr{r}an-e	&	rain	\\
	&	*qa\tbr{z}ay	&	a\tbr{r}a-i	&	a\tbr{r}a	&	jaw, chin	\\	\midrule
\tbx{10-45}	&	*\tbr{d}uha	&	\tbr{r}ua	&	\tbr{r}ua	&	two	\\
	&	*\tbr{d}uRi	&	\tbr{r}uri-n	&	\tbr{r}uri	&	bone (*R > \it{r} through progressive assimilation)	\\
	&	*\tbr{d}aRaq	&	\tbr{r}ara-ne	&	\tbr{r}ara	&	blood	\\
	&	*ma-u\tbr{d}əhi	&	mu\tbr{r}i	&	ana mu{\tl}mu\tbr{r}i	&	\ili{Yalahatan}: `back, behind', \ili{Atamanu}: `youngest child'	\\
	&	*kə\tbr{d}əŋ	&	e\tbr{r}e	&		&	stand	\\	\midrule
\tbx{10-46}	&	*\tbr{d}ahun	&	\tbr{l}ain {\tl} laʔin	&	ai \tbr{r}ain	&	leaf	\\
	&	*\tbr{d}aləm	&	\tbr{l}ale-i	&	pali\tbr{r}are	&	inside	\\
%	&	*tuqə\tbr{d}	&	a-ta{\tl}tue\tbr{l}-e{\footnotemark} 	&		&	tree stump supporting something (tree stump as seat; tree stump supporting raised ritual houses)	\\
	&	*tuqə\tbr{d}	&	a-ta{\tl}tue\tbr{l}-e{\footnotemark} 	&		&	tree stump supporting something	\\
\lspbottomrule
\end{tabularx}\mbox{}\\}
\footnotetext{It is unclear if the /l/ in \it{a-ta{\tl}tuele}
		is a reflex of *d > \it{l},
		or the instrumental suffix \it{-l} (compare \it{aheule} ‘wooden
		comb’, \it{sisiule} ‘wooden comb’, \it{ailoile} ‘sewing needle’; also
		found in other languages; see \citealp[518]{GrimesGrimes2020}).}

The different reflexes of *d in \ili{Yalahatan} in \qf{10-45}
and \qf{10-46} are difficult to explain.
While all cases of *d > \it{l} occur before /{\#}{\gap}a or /{\gap}{\#},
*daRaq > \it{rara-ne} ‘blood’ shows *d > \it{r} in this same environment.
\ili{Yalahatan} \it{lale-i} ‘inside’ could be explained as progressive
assimilation, while \ili{Atamanu} \it{palirare} ‘inside’ shows *l > \it{r} through regressive assimilation.
Given the consistency of \ili{Atamanu} *d/*z > \it{r} it appears the
best solution at this point is to posit *d/*z > **r with subsequent
**r > \it{l, r} in \ili{Yalahatan}.

To summarise, while \ili{Yalahatan}
and \ili{Atamanu} show a merger of *R/*l/*j > \it{l} in most cases,
the reflexes of *d/*z are distinct as **r,
as are the reflexes of *j word finally.
This is similar to \tsc{Patakai-Manusela} with distinct reflexes
of *d/*z compared with *R/*l/*j.